% ============================================================
%  Appendix F — Dataset details for HistoriQA-ThirdRepublic
%  Detailed corpus statistics, historical context, and question
%  typology moved out of the main paper (Sec.~\ref{sec:dataset}).
% ============================================================

\paragraph{Historical context.}
The corpus covers 1887, a dense political year (Boulangist crisis,
Franco-German tensions, parliamentary scandals) in which parliamentary
debates and the contemporaneous \textit{grande presse} were deeply
intertwined: high-circulation dailies acted as both amplifiers of and
partisan participants in political life. This makes the corpus a natural
setting for multi-hop reasoning across institutional discourse and its
mediation in the public sphere, with radically different rhetorical
registers and political commitments.

\paragraph{Corpora.}
Three OCR collections obtained from Gallica, the Biblioth{\`e}que
nationale de France digital library\footnote{\url{https://gallica.bnf.fr}}:
\textit{Les Débats parlementaires} (parliamentary transcripts of the
Chambre des Députés) and two ideologically opposed newspapers,
\textit{Le Gaulois} (monarchist-conservative) and \textit{L'Intransigeant}
(socialist-to-populist). Each collection is indexed as an independent
dense vector store (Appendix~\ref{sec:appendix-experimental-setup}). Key
statistics are in Table~\ref{tab:corpus-stats}.

\begin{table}[htbp]
\centering
\small
\resizebox{\columnwidth}{!}{%
\begin{tabular}{lcccc}
\toprule
 & \multicolumn{2}{c}{\textit{Les Débats}} & \multirow{2}{*}{\makecell{\textit{Le}\\\textit{Gaulois}}} & \multirow{2}{*}{\makecell{\textit{L'Intr-}\\\textit{ansigeant}}} \\
\cmidrule(lr){2-3}
 & All & Filtered & & \\
\midrule
\# documents      & 3,229 & 963   & 78   & 79   \\
Avg tokens/doc    & 2,020 & 4,433 & 922  & 376  \\
Median tokens/doc &   279 &   716 & 802  & 317  \\
Min--max tokens   & 1--63k & 143--63k & 51--4k & 6--1k \\
\bottomrule
\end{tabular}
}
\caption{Corpus statistics. \textit{Les Débats} is reported both raw and
after filtering out non-debate sections (agendas, vote lists).}
\label{tab:corpus-stats}
\end{table}

\paragraph{Multi-hop benchmark overview.}
We evaluate exclusively on the multi-hop subset of
HistoriQA-ThirdRepublic~\cite{historiqa-pellet}: 885 French, historian-validated
questions that each require evidence from at least two distinct
documents. These questions decompose along two orthogonal axes:
a \emph{source-pair} axis (571 newspaper~$\to$~\textit{Les Débats},
314 cross-newspaper) and a \emph{category} axis
(Follow-up, Bridge-Entity, Comparative).

\paragraph{Candidate pairing.}
Multi-hop generation starts from a candidate-pairing stage that selects
document pairs likely to support a single coherent question. All
chunks are embedded with Cohere \texttt{embed-v4.0} and indexed in
independent ChromaDB collections (one per corpus); pairs are scored by
cosine similarity and retained above a fixed threshold $\tau{=}0.7$
calibrated on a held-out historian-annotated sample. Two pairing modes
are used:
\begin{itemize}[leftmargin=1.2em,itemsep=0.15em,topsep=0.25em]
  \item \textit{Newspaper $\to$ Les Débats} (571 questions). For each
    newspaper chunk, the top-1 most similar parliamentary passage is
    paired with it; this direction is preferred over the reverse
    (Débats $\to$ newspaper) to bound the candidate count and maximise
    the diversity of reasoning chains.
  \item \textit{Cross-newspaper} (Gaulois $\leftrightarrow$
    L'Intransigeant, 314 questions). All inter-newspaper pairs above
    $\tau$ are kept, with an additional temporal constraint that
    publication dates differ by at most seven days, ensuring the two
    articles discuss the same news cycle.
\end{itemize}
Pairs purely internal to \textit{Les Débats} were excluded: multi-hop
reasoning over adjacent or identical sessions yields limited
perspective diversity. Pairing across the press/parliament boundary
instead exposes the interpretive gap between official transcripts and
their mediation by partisan dailies.

\paragraph{Historian-in-the-loop generation.}
For each retained pair, a generation prompt is issued to Cohere
\texttt{command-a}. Prompts were iteratively refined with the
domain historian on our team and define three question categories,
inspired by HotpotQA~\cite{hotpotqa_2018}:
\begin{itemize}[leftmargin=1.2em,itemsep=0.15em,topsep=0.25em]
  \item \textit{Follow-up}: an opinion or fact stated in one source and
    the reaction it elicits in the other (e.g., a debate intervention
    and its press commentary).
  \item \textit{Bridge-Entity}: a shared referent (person, institution,
    event) that licenses a chain from one passage to the other.
  \item \textit{Comparative}: contrast between the positions defended in
    each source.
\end{itemize}
The model returns a structured JSON object per question with fields
\texttt{question}, \texttt{answer},
\texttt{supporting\_passages.\{debate,journal\}},
\texttt{reasoning\_type} (\textit{factuel, chronologique, causal,
comparatif, critique/inférentiel}) and \texttt{difficulty}
(\textit{easy/medium/hard}), or the sentinel string
\texttt{AUCUNE QUESTION} when the pair does not license a genuine
multi-hop question. Prompts forbid vague back-references
(\textit{``selon l'article''}) and enforce that questions be
self-contained (explicit date, topic, and main actors).

\paragraph{Validation and release.}
All retained questions were reviewed by the team historian; questions
relying on hallucinated content or trivially answerable from a single
source were discarded. The full system prompts (five variants, one per
category $\times$ source-pair combination), few-shot examples, and the
filtering pipeline are documented in the companion data
paper~\cite{historiqa-pellet} and released in the accompanying
repository.\footnote{\url{https://anonymous.4open.science/r/ORDER_EMNLP2026-BEC6/}}
